% ---- pg-286 （书页 274） ----

\noindent (c) 利用 (a) 的结论, 能否得出如下的结论: 地块(或品种)效应是可加的? 如不可加, 在哪些变量上有交互效应存在, 哪些不存在? 通过作三个独立的双因子一元 ANOVA 来检验.

\noindent (d) 较大的数字对应于高的产量和较好的等级指标. 在地块 2 中, 是否能得出这样的结论: 有一个品种的每个指标均优于其余两个品种? 利用对各对品种的 95\% 庞弗罗尼联合置信区间说明你的结论.

\begin{mvtable}
\phantomsection\label{tab:6.17}\mvtabcap{表 6.17\quad 花生数据}
\begin{tabular}{ccrrr}
\toprule
因子 1 & 因子 2 & $x_1$ & $x_2$ & $x_3$\\
地块 & 品种 & 产量 & 果仁成熟度 & 种子大小\\\midrule
1 & 5 & 195.3 & 153.1 & 51.4\\
1 & 5 & 194.3 & 167.7 & 53.7\\
2 & 5 & 189.7 & 139.5 & 55.5\\
2 & 5 & 180.4 & 121.1 & 44.4\\
1 & 6 & 203.0 & 156.8 & 49.8\\
1 & 6 & 195.9 & 166.0 & 45.8\\
2 & 6 & 202.7 & 166.1 & 60.4\\
2 & 6 & 197.6 & 161.8 & 54.1\\
1 & 8 & 193.5 & 164.5 & 57.8\\
1 & 8 & 187.0 & 165.1 & 58.6\\
2 & 8 & 201.5 & 166.8 & 65.0\\
2 & 8 & 200.0 & 173.8 & 67.2\\
\bottomrule
\end{tabular}
\end{mvtable}

\vspace{2pt}
\noindent\rule{\linewidth}{0.4pt}

\noindent\small 数据来源: 数据经 Yolanda Lopez 许可使用.

\vspace{1em}
\noindent 6.32\quad 在一个试验中, 对处于成长季节的三个品种的一岁树苗观测其不同波长下的遥感和光谱反射数据. 树苗在两种不同水平的营养下成长: 最优水平(用+号表示)和次最优水平(用-号表示). 所用的三种树苗分别是: 云杉(SS), 日本落叶松(JL), 黑松(LP). 两个观测变量是:

\noindent $X_1=$ 在波长为 560 纳米处光谱反射百分比(绿色)

\noindent $X_2=$ 在波长为 720 纳米处光谱反射百分比(红外线附近)

在 235 个阳历日中, 测得在品种和营养水平的不同组合下的相应的子房室均值(CM)数据如下表. 这些平均值均基于四次重复观测.

\begin{mvtable}
\begin{tabular}{rrcc}
\toprule
560CM & 720CM & 品种 & 营养水平\\\midrule
10.35 & 25.93 & SS & +\\
13.41 & 38.63 & JL & +\\
7.78 & 25.15 & LP & +\\
10.40 & 24.25 & SS & -\\
17.78 & 41.45 & JL & -\\
10.40 & 29.20 & LP & -\\
\bottomrule
\end{tabular}
\end{mvtable}

\noindent (a) 将子房室均值看做单个观测, 进行双因子 MANOVA 分析来检验品种效应和营养效应作用. 取 $\alpha=0.05$.

\noindent (b) 分别对 560CM 观测和 720CM 观测构造双因子 ANOVA 表. 这些结果和 (a) 中的
